
cas où la base est un trait. Il s'ensuit que les propriétés sur l'homomorphisme de spécialisation passées en revue dans les sections précédentes, dans le cas où la base est un trait, s'étendent automatiquement au cas d'une base plus générale. Bien entendu, les procédés de définition envisagés s'appliquent tout aussi bien à @@M0@@, @@M1@@ (dans le cas où @@M2@@ est lisse) et @@M3@@ (quand de plus @@M4@@ est séparé).

\textbf{7.17.7.} Pour terminer, signalons que la définition d'un homomorphisme (7.17.3.1) ou (7.17.3.2), ou celle d'un homomorphisme (7.17.4.2), dépend de façon essentielle de certains choix arbitraires (un trait convenable au-dessus de @@M5@@ dans 7.17.3, une suite de @@M6@@ dans 7.17.4, et de plus la donnée d'une résolution des singularités dans 7.17.5), et de même pour les variantes pour @@M7@@ etc. On construit des exemples évidents, avec pour @@M8@@ un point double ordinaire d'une courbe algébrique disons, et @@M9@@ un schéma projectif et lisse sur @@M10@@ (par exemple un schéma abélien de dimension relative 2), où l'homomorphisme de spécialisation sur @@M11@@ et même sur @@M12@@ (groupe de Néron--Sévéri) dépend effectivement du choix du trait @@M13@@ choisi au-dessus de @@M14@@, ou ce qui revient au même ici, du choix d'un point au-dessus de @@M15@@ du normalisé de @@M16@@.

\subsection*{7.18. Appendice à l'Appendice : extension du corps de base pour certains foncteurs en groupes commutatifs}
Nous explicitons ici un point technique soulevé dans 7.17.4.

\textbf{7.18.1.} Soient @@M17@@ un corps, @@M18@@ la catégorie des @@M19@@-algèbres intègres et des homomorphismes plats de telles algèbres, @@M20@@ un foncteur sur @@M21@@, à valeurs dans la catégorie @@M22@@ des ensembles. Supposons que @@M23@@ commute aux limites inductives filtrantes. Soit @@M24@@ une extension de @@M25@@, nous cherchons des cri%
